\documentclass[10pt,a4paper,twoside,bg=print,twocolumn,openany,nodeprecatedcode]{dndbook}
\usepackage[utf8]{inputenc}
\usepackage{tabulary}
\usepackage[normalem]{ulem}
\usepackage{color,soul}
\usepackage{float}
\usepackage{caption}
\usepackage{wrapfig}
\usepackage{tocloft}
\usepackage[hidelinks]{hyperref}

\raggedbottom

\begin{document}
\frontmatter
\rpgMakeCover[
    image = images/adventure/VNotEE/000-00-001.something-wicked-approaches.png,
    title = Vecna: Nest of the Eldritch Eye,
    subtitle = An offsite backup of the preview adventure just for Ben.,
]

\mainmatter


\chapter*{Vecna: Nest of the Eldritch Eye}
\DndDropCapLine{V}{ecna:} Nest of the Eldritch Eye is a fifth edition \textsc{Dungeons \& Dragons} adventure for four to six characters of 3rd level. This adventure takes place in the city of Neverwinter on the Sword Coast. A sinister cult has crept into the city's bowels, and the characters must infiltrate the cult's hideout and root out its members before harm befalls the city.

The information presented here is for the DM's eyes only. If you're planning to play through the adventure with someone else as your DM, stop reading now!


\section{Running the Adventure}
This adventure is designed to be playable in two or three hours. The characters are contracted by Lord Protector Dagult Neverember of Neverwinter to patrol the city. When a dead man is found clutching a withered eyeball, the characters are tasked with uncovering the cause of the man's death. This leads the characters to investigate the city's old catacombs and confront cultists of the lich-god Vecna.
\begin{DndReadAloud}{}
Text that appears in a box like this is meant to be read aloud or paraphrased for the players when their characters first arrive at a location or under a specific circumstance, as described in the text.
\end{DndReadAloud}

When a creature's name appears in bold type, that's a visual clue pointing you to its stat block as a way of saying, "Hey, DM, get this creature's stat block ready. You're going to need it." All monster stat blocks referenced in this adventure can be found in the Monster Manual.

You can make the adventure easier or harder, or adjust it for smaller or larger groups of player characters, by adjusting the number of monsters or by adding or removing encounters.
\section{Beginning the Adventure}
Begin the adventure by reading aloud the following:
\begin{DndReadAloud}{}
The city of Neverwinter stands proudly on the Sword Coast, west and north of Dessarin Valley. Fifty years ago, the city was nearly destroyed by the eruption of Mount Hotenow. Today, the city stands mostly rebuilt. It bustles with skilled tradespeople, intrepid adventurers, and hardy townsfolk.

Leadership in Neverwinter falls to Dagult Neverember, lord protector of the city—and your employer. The city lacks a formal militia, so Lord Neverember often hires mercenaries and adventurers such as yourselves to keep the city secure.
\end{DndReadAloud}

Regardless of whether the characters know each other prior to their arrival in Neverwinter, they are assigned to patrol the city together. The first few days of the patrol are uneventful (unless you decide otherwise), allowing the characters to familiarize themselves with the city.
\section{Dead Man's Message}
After the characters' first few days patrolling the city, tragedy strikes. Read or paraphrase the following:
\begin{DndReadAloud}{}
An anguished scream erupts from a nearby alley. The sound comes from a slender human woman clad in the light leather armor of a scout or investigator. Clutched in her arms is the dead body of a human man clad in a tattered gray robe. The man and woman bear a striking resemblance to each other.

With a dull thump, a desiccated eyeball rolls out of the dead man's palm.
\end{DndReadAloud}

The dead man is Delvin Fearnehart, an investigator in Lord Neverember's employ; the woman holding his body is Kevori Fearnehart, Delvin's sister. A character who succeeds on a DC 8 Intelligence (History) check recognizes both Delvin and Kevori as mercenaries employed by Lord Neverember to patrol the city.

If the characters question Kevori (neutral good, human scout), she shares the following information:

\begin{description}
  \item[Investigators for Hire.]
  Kevori and her brother were tasked by Lord Neverember to investigate suspicious cult activity in the city.
  \item[Rendezvous Gone Wrong.]
  Kevori came here to rendezvous with Delvin after he'd infiltrated a potential cult hideout.
  \item[What She Knows.]
  The only clue Kevori has is the desiccated eyeball that dropped from Delvin's hand.
\end{description}

Kevori requests the characters' aid in uncovering who or what killed her brother while she returns to the Hall of Justice to inform Lord Neverember of these developments. She implores the characters to avenge her brother's death, but she also reminds the characters that Lord Neverember values information and will pay more for capturing cultists alive.

\subsection{Investigating the Eyeball}
The eyeball serves as both a map and a key to the cult's hideout. When a creature holding the eyeball speaks the key phrase, "Hail the Undying," the eyeball glows with green light before spinning to look in the direction of the cult's hideout (acting like a compass).

A Detect Magic spell or similar effect reveals an aura of divination magic around the eyeball. Casting Identify while targeting the eyeball reveals its function and its key phrase. Alternatively, a character can rifle through Delvin's pockets and find a slip of parchment with the key phrase scribbled on it.

Following the activated eyeball's directions leads the characters to a catacomb entrance in one of the still-ruined sections of Neverwinter (see "Into Neverwinter's Catacombs").

\section{Into Neverwinter's Catacombs}
Several months ago, cultists dedicated to Vecna commandeered sections of Neverwinter's catacombs for their nefarious purposes. The sect the characters are pursuing is in the western ruins of the city.

Unless you decide otherwise, the characters encounter no difficulties following the eyeball's directions to a dilapidated entrance into the city's western catacombs. Read or paraphrase the following:
\begin{DndReadAloud}{}
The withered eyeball aims its empty gaze down an alley strewn with rubble. Rats scurry between chunks of old stone and dusty crates. At the back of the alley is a set of moldering wood boards, propped up just enough to block a dark tunnel leading below the surface.
\end{DndReadAloud}

The boards blocking the entrance to the catacombs can be moved easily, revealing a sloping tunnel (area N1).
\subsection{Area Features}

This adventure explores only a small section of Neverwinter's sprawling catacombs. Unless otherwise stated, areas of the catacombs shown in this adventure have the following features:

\begin{description}
\item[Ceilings.]
The ceilings throughout the catacombs are 10 feet high and made of roughly hewn stone.
\item[Eyeball Compass.]
The activated eyeball continuously points in the direction of the cult hideout's entrance (area N11). The eyeball is also needed to unlock the door in area N11.
\item[Light.]
Rusted wall sconces hold lit torches. The area is dimly lit.
\item[Walls and Floors.]
Walls are constructed of stone blocks and are slick with moisture, making climbing impossible without magic or gear. Floors are made from a mix of stone slabs and packed earth.
\end{description}

\subsection{Area Locations}
The following locations are keyed to the Nest of the Eldritch Eye map.

\begin{figure*}
\includegraphics[width=\textwidth]{images/adventure/VNotEE/001-map-0.01-catacombs.png}
\end{figure*}

\subsubsection{N1: Entrance Tunnel}
The tunnel proceeds at a downward slope, bringing the characters 20 feet beneath street level.

\subsubsection{N2: Half-Collapsed Chamber}

\begin{DndReadAloud}{}
The soft lapping of water echoes through this wide chamber. Wood and stone debris from collapsed buildings litter the area. Brackish water floods nearly half of the chamber, getting deeper toward a dark tunnel to the west.
\end{DndReadAloud}

Characters who have a passive Wisdom (Perception) score of 12 or higher notice the entrance to a cramped tunnel hidden behind some collapsed rubble to the south. A character who spends 1 minute clearing the rubble opens the entrance to a crawlspace (area N7).

\subparagraph{Making a Skiff}
Characters have the option of assembling a makeshift skiff from the debris here to make traversing the flooded area easier. A character attempting to assemble a skiff must make a DC 11 Wisdom (Survival) check. On a successful check, the character makes a skiff sturdy enough to support up to two Medium or smaller creatures.

While riding a skiff, a creature can use its movement to propel the skiff up to 20 feet in a direction of its choice.

\subsubsection{N3: Flooded Tunnel}
The water filling this tunnel is 3 feet deep; without the use of a skiff or other water vehicle, the water counts as difficult terrain for Medium or larger creatures that lack a swimming speed, while Small and smaller creatures must swim to traverse the tunnel.

\subparagraph{Water Weird}
A neutral evil water weird lurks here. It tries to ambush and kill any creature that enters the tunnel.

Upon entering the tunnel for the first time, a character who has a passive Wisdom (Perception) score of 15 or higher notices odd waves and currents that suggest the presence of a creature in the water. If a character's passive Wisdom (Perception) score is lower than 15, the character has disadvantage on their initiative roll when the water weird attacks.

\subsubsection{N4: Cinerary Rotunda}

\begin{DndReadAloud}{}
Glossy urns and cinerary boxes line the walls of this vaulted rotunda. Drifting in the center of the room above a drain is a ghostly, transparent, humanoid figure. The figure is clad in plate armor, and where a head should be, there is instead a featureless, luminescent orb.
\end{DndReadAloud}

The neutral ghost haunting this area has no memory of its name or existence in life—a fact which pains it greatly. Upon noticing the characters enter, the ghost beseeches them for help in recovering its previous identity. If the characters refuse, the ghost sulks but isn't hostile.
\subparagraph{Discovering the Ghost's Identity}
The ghost is the spirit of the knight entombed in area N6. Clues regarding the ghost's identity can be found in areas N5 and N6. A character who searches the rotunda for 10 minutes finds no clues related to the ghost in this room.

To regain its identity, the ghost must hear its full name: Chanelle Hallwinter. If a character speaks this name aloud within 20 feet of the ghost, the featureless orb of the ghost's head coalesces into the face of a middle-aged human woman with braided hair, and the ghost's memories return. The ghost is grateful to the characters and resolves to aid them.
\subparagraph{What the Ghost Knows}

At first, the ghost remembers nothing about Delvin, the cult, or the catacombs. Once the ghost's identity has been restored, it remembers not only its life as Chanelle Hallwinter but also what it has witnessed while haunting these catacombs. The ghost can share the following information with the characters:

\begin{description}
\item[Catacomb Invaders.]
A small group of Humanoids descended into the catacombs several months ago, whispering about strange rituals and secrets.
\item[Delvin.]
Chanelle saw a man matching Delvin's description flee through the rotunda from deeper within the catacombs, pursued by an individual in a hooded gray robe.
\item[Zombie Infestation.]
The crypt to the south of the rotunda was desecrated and now swarms with zombies.
\end{description}

\subsubsection{N5: Tomb of the Mage}
\begin{DndReadAloud}{}
The smell of old parchment fills this tomb. Bookshelves carrying various tomes and scholarly implements stand against the walls. At the center is a stone, gold-painted sarcophagus, the sides of which bear a beautiful relief carving of two humanoid women exploring a forest together.
\end{DndReadAloud}

This is the tomb of Makalia Siannodel, a female elf mage and the partner of the knight entombed in area N6.

\subparagraph{Discovering the Ghost's Identity}
Characters who search this tomb find treasure (see "Treasure" below) and the following clues:

\begin{description}
\item[Golden Locket.]
A heart-shaped golden locket sits on one of the bookshelves. The locket is locked but can be either pried open with a successful DC 11 Strength (Athletics) check or picked open with a successful DC 10 Dexterity (Sleight of Hand) check using thieves' tools. Inside is a small portrait of a human woman, next to which is the inscription, "My dearest Chanelle."
\item[Relief Details.]
One of the figures depicted in the sarcophagus's relief carving wears armor that matches the armor of the ghost in area N4.
\end{description}

\subparagraph{Treasure}
Nestled among the bookshelves are five incense sticks (worth 10 gp each) and a component pouch. A character who inspects the area and succeeds on a DC 15 Intelligence (Investigation) check also finds a hidden compartment in one of the shelved books. Inside is a Spell Scroll of Protection from Evil and Good.
\subsubsection{N6: Tomb of the Knight}
\begin{DndReadAloud}{}
A stone, gold-painted sarcophagus rests in the center of this tomb. The sides of the sarcophagus bear a relief depicting two women gazing lovingly over a field and a city, and the lid has a faintly distinguishable family crest carved into it. Old weapons and decorated armor line the walls of this room.
\end{DndReadAloud}

This is the tomb of Chanelle Hallwinter, a female human knight and the partner of the mage entombed in area N5. This is also the tomb of the ghost found in area N4.
\subparagraph{Discovering the Ghost's Identity}
Characters who search this tomb for clues find the following pieces of information:

\begin{description}
\item[Family Crest.]
Engraved on the lid of the sarcophagus is a coat of arms depicting a blunted six-point crown. A character who succeeds on a DC 10 Intelligence (History) check recognizes this symbol as the crest of the Hallwinter family, whose lineage produced renowned knights throughout the Sword Coast.
\item[Relief Details.]
One of the figures depicted in the sarcophagus's relief carving wears armor that matches the armor of the ghost in area N4.
\end{description}

\subparagraph{Treasure}
Most items stored in this tomb have rusted, but one shield remains intact. This shield is a cursed Shield of Missile Attraction; however, if the characters restore the ghost's identity, Chanelle Hallwinter gratefully gives the shield to the characters and removes the curse, transforming the shield into a +1 Shield.
\subparagraph{Secret Passage}
A character who searches the tomb for secret doors and succeeds on a DC 14 Wisdom (Perception) check sees the faint outline of a doorway at the back of the tomb. The door is unlocked, swings open with a gentle push, and leads to area N8.
\subsubsection{N7: Crawlspace}
The cramped passage is large enough to accommodate Small and smaller creatures; Medium creatures must squeeze to move through it.
\subparagraph{Disturbed Rubble}
When a creature enters the crawlspace for the first time, its movement disturbs some of the rubble above. The creature must succeed on a DC 12 Dexterity saving throw to reach the other side of the crawlspace unscathed. On a failed save, the creature takes 7 (3d4) bludgeoning damage from the falling rubble.
\subsubsection{N8: Forked Tunnel}
\begin{DndReadAloud}{}
The tunnel splits three ways: north toward a tight crawlspace, west toward a stone door, and south toward a foul-smelling chamber.
\end{DndReadAloud}

At the end of the west tunnel is a secret door to area N6. From this side, the door is clearly visible and pulls open easily.
\subsubsection{N9: Sewerage Chamber}
\begin{DndReadAloud}{}
A foul stench rises from the murky sewage that fills most of this large chamber. Three four-foot-tall burbling masses slink through the muck.
\end{DndReadAloud}

The sewage is 1 foot deep and difficult terrain.

The three burbling masses are awakened piles of undead sludge, sloughed from the cult's twisted experiments (each uses the gray ooze stat block but is an Undead instead of Ooze). The masses ignore the room's difficult terrain and attack if approached.
\subsubsection{N10: Crypt of the Silenced Singers}
\begin{DndReadAloud}{}
What was once a serene crypt now lies in ruin. Two chipped stone pillars brace the vaulted ceiling. Tombs are cracked open, and reanimated corpses prowl the room. At the center of the room's back wall is a dirtied and desecrated shrine bearing the image of a blank scroll.
\end{DndReadAloud}

Eight Zombie shamble around this crypt. The zombies are hungry and immediately attack upon seeing a creature unless that creature succeeds on a DC 16 Dexterity (Stealth) check.
\subparagraph{Shrine to Oghma}
A character who examines the shrine and succeeds on a DC 10 Intelligence (Religion) check recognizes the image on the shrine as the holy symbol of Oghma, god of knowledge and patron to bards and wizards. If the check succeeds by 3 or more, the character intuits that rededicating the shrine to Oghma could help against the zombies in the area.

To rededicate the shrine, a creature must first use an action to clean it. Once the filth is cleaned off, a creature can use an action to do one of the following:

\begin{itemize}
\item Make a DC 14 Charisma (Performance) check using an instrument within 5 feet of the shrine, performing a song dedicated to Oghma.
\item Touch the shrine and expend a spell slot of 1st level or higher.
\end{itemize}

Once one of the two options has been done successfully, the shrine glows with white light, and all zombies in the area immediately collapse with the unconscious condition for 24 hours. The condition ends early for a zombie if it takes damage.
\subsubsection{N11: Door of the Eldritch Eye}
\begin{DndReadAloud}{}
A double door made of green-tinged stone seals off further progression into the catacombs. Facing the double door is a large carving of a grinning skull; one eye socket bears a wide eyeball with a jeweled iris, while the other is an empty divot.
\end{DndReadAloud}

The double door is locked and has no keyholes. To open the door from the north side, a desiccated eyeball (such as the one obtained from Delvin) must be placed in the divot of the grinning skull. Once placed there, the eyeball pulses with sickly green light, and the double door swings open to reveal a staircase leading further down. If the eyeball is removed from the divot, the door remains open for 1 minute before closing and locking on its own. A Knock spell or similar magic also opens the door.

Neither magic nor an eyeball is needed to open the double door from the south. On the door's south side is a stone pedal; stepping on it causes the double door to swing open.
\subsubsection{N12: Hall of the Whispered One}
\begin{DndReadAloud}{}
The bottom of the stairway opens into a wide sanctuary with a vaulted ceiling. Stone pews are arranged in orderly rows. Black candles burning green flames occupy niches along the walls. Atop a pulpit at the far end of the hall stands a jagged sculpture of an emaciated hand with one eyeball in its palm.
\end{DndReadAloud}

A character who succeeds on a DC 11 Intelligence (Religion) check recognizes the sculpture as the symbol of Vecna, a powerful lich and god of secrets known and feared on many worlds.

Any loud noises here awaken the cultists in area N13, who investigate at once.
\subsubsection{N13: Cultist Quarters}
If the occupants of this room were lured to area N12 by noise there, omit the last sentence of the following boxed text:
\begin{DndReadAloud}{}
Makeshift cots are strewn about two adjoining caverns separated by an open doorway. Several individuals in gray robes slumber here peacefully despite the flickering torchlight.
\end{DndReadAloud}

Seven Cultist sleep in these two chambers. Each cultist wears a hooded gray robe that has a desiccated eyeball (like the one obtained from Delvin) tucked in one pocket.
\subparagraph{Being Sneaky}
To move through these rooms without waking the cultists, a creature must succeed on a DC 14 Dexterity (Stealth) check. On a failed check, 1d4 of the cultists wake up.

Awake cultists are hostile but easily cowed. A character can use an action to make a DC 12 Charisma (Intimidation) check. On a successful check, all awake cultists that can see or hear the character surrender.
\subparagraph{Interrogating the Cultists}
Characters can attempt to glean more information from captured cultists by making a DC 10 Charisma (Intimidation or Persuasion) check. On a successful check, a character learns one of the following pieces of information:

\begin{description}
\item[Cult Leader.]
The cult's leader is Zalryr, who is conducting a profane experiment in the ritual chamber (area N15). Zalryr aims to "draw secrets from the depths of the mortal soul."
\item[Cult Worship.]
The cultists all serve Vecna, whom they refer to as the Whispered One.
\item[Delvin's Fate.]
Zalryr uncovered the presence of a spy recently. That spy was chased from the catacombs and assassinated for his transgression. The cultists' description of the spy confirms it was Delvin Fearnehart.
\end{description}

\subparagraph{Treasure}
A character who searches these rooms and makes a successful DC 12 Intelligence (Investigation) check finds two gray cultist robes and three gaudy jeweled rings worth 10 gp each. If the check succeeds by 3 or more, the character also finds a Hat of Disguise tucked beneath one of the unused cots.
\subsubsection{N14: Eldritch Library}
\begin{DndReadAloud}{}
Crooked shelves filled with books and scrolls stand against the walls. In the center of the room is a square table covered in scribbled notes and ink-stained parchment.
\end{DndReadAloud}

The notes spread across the table detail the cult's experiments and history. A character who spends at least 10 minutes studying the notes learns the following:

\begin{description}
\item[Cult Experiments.]
The cult's leader, Zalryr, is experimenting with ways to magically siphon secrets from an individual's soul. Early experiments reduced volunteers to piles of necrotic sludge, which were disposed of in the sewerage (area N9).

\item[Others.]
The cultists here are but one sect of many that have infiltrated Neverwinter. The other sects have chosen different areas of the city's catacombs as their respective hideouts.

\end{description}

\subparagraph{Trapped Lockbox}
A character who inspects the shelves and succeeds on a DC 10 Intelligence (Investigation) check finds a small ebony box. The box is locked; Zalryr in area N15 has the key. The box can be picked open by a creature that makes a successful DC 13 Dexterity (Sleight of Hand) check using thieves' tools, or it can be forced open with a successful DC 15 Strength (Athletics) check.

The box is trapped. A creature that opens the box by any means other than using the proper key must make a DC 13 Constitution saving throw, taking 7 (2d6) necrotic damage on a failed saved or half as much damage on a successful one. This trap can't be detected or disarmed.

Inside the box are two Potion of Greater Healing.
\subsubsection{N15: Ritual Chamber}
\begin{DndReadAloud}{}
Cacophonous whispers echo through this expansive cave. An imposing human cultist in a sweeping robe stands against the wall, chanting. His oily hair is slicked back, and his skin is gaunt and gray.

Carved into the floor before him is a runic circle that pulses with sickly green light. Two cultists—one human, one elf—stand inside the circle, heads thrown back and mouths agape. Their leader commands, "Now! Release your secrets unto me! Let the truths hidden in your soul come forth and become stronger!" As if in response, the cultists' silhouettes warp into two lanky shadowy entities.
\end{DndReadAloud}

Zalryr (neutral evil, human cult fanatic) stands at the north end of the room, conducting an experiment on the two Cultist standing inside the runic circle. Zalryr's experiment has magically siphoned the secrets from the cultists' souls to create the two shadowy entities (each uses the shadow stat block).

Each cultist, including Zalryr, wears a hooded gray robe that has a desiccated eyeball (like the one obtained from Delvin) tucked in one pocket.

Upon noticing the characters, Zalryr commands the cultists and the shadowy entities to attack. The cultists and entities fight to their deaths, but Zalryr surrenders when he is reduced to 10 hit points or fewer.
\subparagraph{Vecna Appears}
When Zalryr either surrenders or is slain, read or paraphrase the following:
\begin{DndReadAloud}{}
The runic circle suddenly hisses and flashes with lurid light. As if being extinguished, the glow of the runes dims as shadowy smoke rises from the carvings. The smoke gathers in the center of the chamber, where it coalesces into an apparition of an emaciated skull with one glowing green eye.

The apparition speaks in a hissing baritone. "Good news, Zalryr. Though your efforts have been disappointing, you have brought me new... points of interest." The apparition swivels to look at each of you in turn. "Yes, there is great potential here. May we meet again—but until then, I have my eye on you." The apparition then explodes into streaks of shadow, passing through each of you with a whispered scream before vanishing.
\end{DndReadAloud}

Vecna bestows one of two supernatural gifts on each character. Have each player roll a die, with the result determining which charm their character receives: the Charm of the Creeping Hand if the die roll is an odd number or the Charm of the Eldritch Eye if the die roll is an even number. While in possession of either charm, the character has the unshakeable feeling that they are being watched. See the next section for descriptions of these charms.
\subsection{Vecna's Supernatural Gifts}
Vecna is known to bestow supernatural gifts on mortals who impress him, regardless of their affiliations. The following charms are two of his favorites.

\DndItemHeader{Charm of the Creeping Hand}{Charm}
Once per turn, when you hit a creature with an attack roll using a weapon or an unarmed strike, you can infuse your strike with life-stealing energy. Your attack then deals an extra 1d10 necrotic damage, and you gain 5 temporary hit points. Once used five times, the charm vanishes

\DndItemHeader{Charm of the Eldritch Eye}{Charm}
You can cast Clairvoyance as an action, without using a spell slot and requiring no material components. Once used three times, the charm vanishes.


\section{Concluding the Adventure}
The characters can take any captives, leave the catacombs, and report to Lord Neverember without further incident. Lord Neverember is grateful for the characters' efforts and pays each character a base sum of 100 gp for the job, plus an additional 10 gp each for each cultist captured alive (25 gp each for Zalryr).
\subsection{Continuing the Campaign}
Though this sect of cultists has been neutralized, others still lurk in the shadows of Neverwinter, with even more nefarious schemes in service of Vecna. Lord Neverember might call on the characters to continue rooting out these cultists, culminating in an adventure through the sprawling catacombs beneath Neverwinter's Neverdeath Graveyard. If you need help planning further, Vecna: Eve of Ruin contains more information about cult activities in Neverdeath Graveyard and provides an excellent campaign that follows Vecna's machinations.

\appendix

\chapter{Creatures}
\addcontentsline{toc}{section}{Cult Fanatic}

\begin{DndMonster}[float=ht]{Cult Fanatic}
\DndMonsterType{Medium humanoid (any race), any non-good alignment}
\DndMonsterBasics[
armor-class = {13 (leather armor)},
hit-points = {\DndDice{6d8 + 6}},
speed = {30 ft.}
]
\DndMonsterAbilityScores[
 str = 11,
 dex = 14,
 con = 12,
 int = 10,
 wis = 13,
 cha = 14
]
\DndMonsterDetails[
skills = {Deception +4, Persuasion +4, Religion +2},
languages = {any one language (usually Common)},
challenge = 2,
]
\DndMonsterAction{Dark Devotion}
The fanatic has advantage on saving throws against being charmed or frightened.
\DndMonsterAction{Spellcasting}
The fanatic is a 4th-level spellcaster. Its spellcasting ability is Wisdom (spell save DC 11, +3 to hit with spell attacks). The fanatic has the following cleric spells prepared:
\begin{DndMonsterSpells}
  \DndMonsterSpellLevel{light, sacred flame, thaumaturgy}
   \DndMonsterSpellLevel[1][4]{command, inflict wounds, shield of faith}
   \DndMonsterSpellLevel[2][3]{hold person, spiritual weapon}
\end{DndMonsterSpells}
\par
\DndMonsterSection{Actions}
\DndMonsterAction{Multiattack}
The fanatic makes two melee attacks.
\DndMonsterAction{Dagger}
\textsl{Melee or Ranged Weapon Attack:} +4 to hit, reach 5 ft. or range 20/60 ft., one creature. \textsl{Hit:} 4 (1d4 + 2) piercing damage.
\end{DndMonster}


\begin{figure}[h]
\includegraphics[width=0.5\textwidth]{images/bestiary/MM/Cult Fanatic.png}
\end{figure}
\clearpage
\addcontentsline{toc}{section}{Cultist}

\begin{DndMonster}[float=ht]{Cultist}
\DndMonsterType{Medium humanoid (any race), any non-good alignment}
\DndMonsterBasics[
armor-class = {12 (leather armor)},
hit-points = {\DndDice{2d8}},
speed = {30 ft.}
]
\DndMonsterAbilityScores[
 str = 11,
 dex = 12,
 con = 10,
 int = 10,
 wis = 11,
 cha = 10
]
\DndMonsterDetails[
skills = {Deception +2, Religion +2},
languages = {any one language (usually Common)},
challenge = 1/8,
]
\DndMonsterAction{Dark Devotion}
The cultist has advantage on saving throws against being charmed or frightened.
\par
\DndMonsterSection{Actions}
\DndMonsterAction{Scimitar}
\textsl{Melee Weapon Attack:} +3 to hit, reach 5 ft., one creature. \textsl{Hit:} 4 (1d6 + 1) slashing damage.
\end{DndMonster}


\begin{figure}[h]
\includegraphics[width=0.5\textwidth]{images/bestiary/MM/Cultist.png}
\end{figure}
\clearpage
\addcontentsline{toc}{section}{Ghost}

\begin{DndMonster}[float=ht]{Ghost}
\DndMonsterType{Medium undead, any alignment}
\DndMonsterBasics[
armor-class = {11},
hit-points = {\DndDice{10d8}},
speed = {0 ft., fly 40 ft. (hover)}
]
\DndMonsterAbilityScores[
 str = 7,
 dex = 13,
 con = 10,
 int = 10,
 wis = 12,
 cha = 17
]
\DndMonsterDetails[
damage-resistances = {acid, fire, lightning, thunder; bludgeoning, piercing, slashing from nonmagical attacks},
damage-immunities = {cold, necrotic, poison},
senses = {darkvision 60 ft., passive perception 11},
languages = {any languages it knew in life},
challenge = 4,
]
\DndMonsterAction{Ethereal Sight}
The ghost can see 60 feet into the Ethereal Plane when it is on the Material Plane, and vice versa.
\DndMonsterAction{Incorporeal Movement}
The ghost can move through other creatures and objects as if they were difficult terrain. It takes 5 (1d10) force damage if it ends its turn inside an object.
\par
\DndMonsterSection{Actions}
\DndMonsterAction{Withering Touch}
\textsl{Melee Weapon Attack:} +5 to hit, reach 5 ft., one target. \textsl{Hit:} 17 (4d6 + 3) necrotic damage.
\DndMonsterAction{Etherealness}
The ghost enters the Ethereal Plane from the Material Plane, or vice versa. It is visible on the Material Plane while it is in the Border Ethereal, and vice versa, yet it can't affect or be affected by anything on the other plane.
\DndMonsterAction{Horrifying Visage}
Each non-undead creature within 60 feet of the ghost that can see it must succeed on a DC 13 Wisdom saving throw or be frightened for 1 minute. If the save fails by 5 or more, the target also ages 1d4 × 10 years. A frightened target can repeat the saving throw at the end of each of its turns, ending the frightened condition on itself on a success. If a target's saving throw is successful or the effect ends for it, the target is immune to this ghost's Horrifying Visage for the next 24 hours. The aging effect can be reversed with a  greater restoration spell, but only within 24 hours of it occurring.
\DndMonsterAction{Possession Recharge 6}
One humanoid that the ghost can see within 5 feet of it must succeed on a DC 13 Charisma saving throw or be possessed by the ghost; the ghost then disappears, and the target is incapacitated and loses control of its body. The ghost now controls the body but doesn't deprive the target of awareness. The ghost can't be targeted by any attack, spell, or other effect, except ones that turn undead, and it retains its alignment, Intelligence, Wisdom, Charisma, and immunity to being charmed and frightened. It otherwise uses the possessed target's statistics, but doesn't gain access to the target's knowledge, class features, or proficiencies.
The possession lasts until the body drops to 0 hit points, the ghost ends it as a bonus action, or the ghost is turned or forced out by an effect like the dispel evil and good spell. When the possession ends, the ghost reappears in an unoccupied space within 5 feet of the body. The target is immune to this ghost's Possession for 24 hours after succeeding on the saving throw or after the possession ends.
\end{DndMonster}


\begin{figure}[h]
\includegraphics[width=0.5\textwidth]{images/bestiary/MM/Ghost.png}
\end{figure}
\clearpage
\addcontentsline{toc}{section}{Gray Ooze}

\begin{DndMonster}[float=ht]{Gray Ooze}
\DndMonsterType{Medium ooze, unaligned}
\DndMonsterBasics[
armor-class = {8},
hit-points = {\DndDice{3d8 + 9}},
speed = {10 ft., 10 ft. (key)}
]
\DndMonsterAbilityScores[
 str = 12,
 dex = 6,
 con = 16,
 int = 1,
 wis = 6,
 cha = 2
]
\DndMonsterDetails[
skills = {Stealth +2},
damage-resistances = {acid, cold, fire},
senses = {blindsight 60 ft. (blind beyond this radius), passive perception 8},
challenge = 1/2,
]
\DndMonsterAction{Amorphous}
The ooze can move through a space as narrow as 1 inch wide without squeezing.
\DndMonsterAction{Corrode Metal}
Any nonmagical weapon made of metal that hits the ooze corrodes. After dealing damage, the weapon takes a permanent and cumulative −1 penalty to damage rolls. If its penalty drops to −5, the weapon is destroyed. Nonmagical ammunition made of metal that hits the ooze is destroyed after dealing damage.
The ooze can eat through 2-inch-thick, nonmagical metal in 1 round.
\DndMonsterAction{False Appearance}
While the ooze remains motionless, it is indistinguishable from an oily pool or wet rock.
\par
\DndMonsterSection{Actions}
\DndMonsterAction{Pseudopod}
\textsl{Melee Weapon Attack:} +3 to hit, reach 5 ft., one target. \textsl{Hit:} 4 (1d6 + 1) bludgeoning damage plus 7 (2d6) acid damage, and if the target is wearing nonmagical metal armor, its armor is partly corroded and takes a permanent and cumulative −1 penalty to the AC it offers. The armor is destroyed if the penalty reduces its AC to 10.
\end{DndMonster}


\begin{figure}[h]
\includegraphics[width=0.5\textwidth]{images/bestiary/MM/Gray Ooze.png}
\end{figure}
\clearpage
\addcontentsline{toc}{section}{Scout}

\begin{DndMonster}[float=ht]{Scout}
\DndMonsterType{Medium humanoid (any race), any alignment}
\DndMonsterBasics[
armor-class = {13 (leather armor)},
hit-points = {\DndDice{3d8 + 3}},
speed = {30 ft.}
]
\DndMonsterAbilityScores[
 str = 11,
 dex = 14,
 con = 12,
 int = 11,
 wis = 13,
 cha = 11
]
\DndMonsterDetails[
skills = {Nature +4, Perception +5, Stealth +6, Survival +5},
languages = {any one language (usually Common)},
challenge = 1/2,
]
\DndMonsterAction{Keen Hearing and Sight}
The scout has advantage on Wisdom (Perception) checks that rely on hearing or sight.
\par
\DndMonsterSection{Actions}
\DndMonsterAction{Multiattack}
The scout makes two melee attacks or two ranged attacks.
\DndMonsterAction{Shortsword}
\textsl{Melee Weapon Attack:} +4 to hit, reach 5 ft., one target. \textsl{Hit:} 5 (1d6 + 2) piercing damage.
\DndMonsterAction{Longbow}
\textsl{Ranged Weapon Attack:} +4 to hit, ranged 150/600 ft., one target. \textsl{Hit:} 6 (1d8 + 2) piercing damage.
\end{DndMonster}


\begin{figure}[h]
\includegraphics[width=0.5\textwidth]{images/bestiary/MM/Scout.png}
\end{figure}
\clearpage
\addcontentsline{toc}{section}{Shadow}

\begin{DndMonster}[float=ht]{Shadow}
\DndMonsterType{Medium undead, chaotic evil}
\DndMonsterBasics[
armor-class = {12},
hit-points = {\DndDice{3d8 + 3}},
speed = {40 ft.}
]
\DndMonsterAbilityScores[
 str = 6,
 dex = 14,
 con = 13,
 int = 6,
 wis = 10,
 cha = 8
]
\DndMonsterDetails[
skills = {Stealth +4},
damage-vulnerabilities = {radiant},
damage-resistances = {acid, cold, fire, lightning, thunder; bludgeoning, piercing, slashing from nonmagical attacks},
damage-immunities = {necrotic, poison},
senses = {darkvision 60 ft., passive perception 10},
challenge = 1/2,
]
\DndMonsterAction{Amorphous}
The shadow can move through a space as narrow as 1 inch wide without squeezing.
\DndMonsterAction{Shadow Stealth}
While in dim light or darkness, the shadow can take the Hide action as a bonus action. Its stealth bonus is also improved to +6.
\DndMonsterAction{Sunlight Weakness}
While in sunlight, the shadow has disadvantage on attack rolls, ability checks, and saving throws.
\par
\DndMonsterSection{Actions}
\DndMonsterAction{Strength Drain}
\textsl{Melee Weapon Attack:} +4 to hit, reach 5 ft., one creature. \textsl{Hit:} 9 (2d6 + 2) necrotic damage, and the target's Strength score is reduced by 1d4. The target dies if this reduces its Strength to 0. Otherwise, the reduction lasts until the target finishes a short or long rest.
If a non-evil humanoid dies from this attack, a new shadow rises from the corpse 1d4 hours later.
\end{DndMonster}


\begin{figure}[h]
\includegraphics[width=0.5\textwidth]{images/bestiary/MM/Shadow.png}
\end{figure}
\clearpage
\addcontentsline{toc}{section}{Water Weird}

\begin{DndMonster}[float=ht]{Water Weird}
\DndMonsterType{Large elemental, neutral}
\DndMonsterBasics[
armor-class = {13},
hit-points = {\DndDice{9d10 + 9}},
speed = {0 ft., 60 ft. (key)}
]
\DndMonsterAbilityScores[
 str = 17,
 dex = 16,
 con = 13,
 int = 11,
 wis = 10,
 cha = 10
]
\DndMonsterDetails[
damage-resistances = {fire; bludgeoning, piercing, slashing from nonmagical attacks},
damage-immunities = {poison},
senses = {blindsight 30 ft., passive perception 10},
languages = {understands Aquan but doesn't speak},
challenge = 3,
]
\DndMonsterAction{Invisible in Water}
The water weird is invisible while fully immersed in water.
\DndMonsterAction{Water Bound}
The water weird dies if it leaves the water to which it is bound or if that water is destroyed.
\par
\DndMonsterSection{Actions}
\DndMonsterAction{Constrict}
\textsl{Melee Weapon Attack:} +5 to hit, reach 10 ft., one creature. \textsl{Hit:} 13 (3d6 + 3) bludgeoning damage. If the target is Medium or smaller, it is grappled (escape DC 13) and pulled 5 feet toward the water weird. Until this grapple ends, the target is restrained, the water weird tries to drown it, and the water weird can't constrict another target.
\end{DndMonster}


\begin{figure}[h]
\includegraphics[width=0.5\textwidth]{images/bestiary/MM/Water Weird.png}
\end{figure}
\clearpage
\addcontentsline{toc}{section}{Zombie}

\begin{DndMonster}[float=ht]{Zombie}
\DndMonsterType{Medium undead, neutral evil}
\DndMonsterBasics[
armor-class = {8},
hit-points = {\DndDice{3d8 + 9}},
speed = {20 ft.}
]
\DndMonsterAbilityScores[
 str = 13,
 dex = 6,
 con = 16,
 int = 3,
 wis = 6,
 cha = 5
]
\DndMonsterDetails[
saving-throws = {Wis +0},
damage-immunities = {poison},
senses = {darkvision 60 ft., passive perception 8},
languages = {understands all languages it spoke in life but can't speak},
challenge = 1/4,
]
\DndMonsterAction{Undead Fortitude}
If damage reduces the zombie to 0 hit points, it must make a Constitution saving throw with a DC of 5 + the damage taken, unless the damage is radiant or from a critical hit. On a success, the zombie drops to 1 hit point instead.
\par
\DndMonsterSection{Actions}
\DndMonsterAction{Slam}
\textsl{Melee Weapon Attack:} +3 to hit, reach 5 ft., one target. \textsl{Hit:} 4 (1d6 + 1) bludgeoning damage.
\end{DndMonster}


\begin{figure}[h]
\includegraphics[width=0.5\textwidth]{images/bestiary/MM/Zombie.png}
\end{figure}
\clearpage
\chapter{Spells}
\addcontentsline{toc}{section}{Detect Magic}


\DndSpellHeader%
{Detect Magic}
{1st level Divination}
{1 action}
{Self}
{V, S}
{Concentration, up to 10 minutes}
For the duration, you sense the presence of magic within 30 feet of you. If you sense magic in this way, you can use your action to see a faint aura around any visible creature or object in the area that bears magic, and you learn its school of magic, if any.
The spell can penetrate most barriers, but it is blocked by 1 foot of stone, 1 inch of common metal, a thin sheet of lead, or 3 feet of wood or dirt.
\addcontentsline{toc}{section}{Dispel Evil and Good}


\DndSpellHeader%
{Dispel Evil and Good}
{5th level Abjuration}
{1 action}
{Self}
{V, S, holy water or powdered silver and iron}
{Concentration, up to 1 minute}
Shimmering energy surrounds and protects you from fey, undead, and creatures originating from beyond the Material Plane. For the duration, celestials, elementals, fey, fiends, and undead have disadvantage on attack rolls against you.
You can end the spell early by using either of the following special functions.
\subparagraph{Break Enchantment}
As your action, you touch a creature you can reach that is charmed, frightened, or possessed by a celestial, an elemental, a fey, a fiend, or an undead. The creature you touch is no longer charmed, frightened, or possessed by such creatures.
\subparagraph{Dismissal}
As your action, make a melee spell attack against a celestial, an elemental, a fey, a fiend, or an undead you can reach. On a hit, you attempt to drive the creature back to its home plane. The creature must succeed on a Charisma saving throw or be sent back to its home plane (if it isn't there already). If they aren't on their home plane, undead are sent to the Shadowfell, and fey are sent to the Feywild.
\addcontentsline{toc}{section}{Greater Restoration}


\DndSpellHeader%
{Greater Restoration}
{5th level Abjuration}
{1 action}
{Touch}
{V, S, diamond dust worth at least 100 gp, which the spell consumes}
{Instantaneous}
You imbue a creature you touch with positive energy to undo a debilitating effect. You can reduce the target's exhaustion level by one, or end one of the following effects on the target:

\begin{itemize}
\item One effect that charmed or petrified the target
\item One curse, including the target's attunement to a cursed magic item
\item Any reduction to one of the target's ability scores
\item One effect reducing the target's hit point maximum
\end{itemize}

\addcontentsline{toc}{section}{Identify}


\DndSpellHeader%
{Identify}
{1st level Divination}
{1 minute}
{Touch}
{V, S, a pearl worth at least 100 gp and an owl feather}
{Instantaneous}
You choose one object that you must touch throughout the casting of the spell. If it is a magic item or some other magic-imbued object, you learn its properties and how to use them, whether it requires attunement to use, and how many charges it has, if any. You learn whether any spells are affecting the item and what they are. If the item was created by a spell, you learn which spell created it.
If you instead touch a creature throughout the casting, you learn what spells, if any, are currently affecting it.
\addcontentsline{toc}{section}{Knock}


\DndSpellHeader%
{Knock}
{2nd level Transmutation}
{1 action}
{60 feet}
{V}
{Instantaneous}
Choose an object that you can see within range. The object can be a door, a box, a chest, a set of manacles, a padlock, or another object that contains a mundane or magical means that prevents access.
A target that is held shut by a mundane lock or that is stuck or barred becomes unlocked, unstuck, or unbarred. If the object has multiple locks, only one of them is unlocked.
If you choose a target that is held shut with arcane lock, that spell is suppressed for 10 minutes, during which time the target can be opened and shut normally.
When you cast the spell, a loud knock, audible from as far away as 300 feet, emanates from the target object.
\addcontentsline{toc}{section}{Protection from Evil and Good}


\DndSpellHeader%
{Protection from Evil and Good}
{1st level Abjuration}
{1 action}
{Touch}
{V, S, holy water or powdered silver and iron, which the spell consumes}
{Concentration, up to 10 minutes}
Until the spell ends, one willing creature you touch is protected against certain types of creatures: aberrations, celestials, elementals, fey, fiends, and undead.
The protection grants several benefits. Creatures of those types have disadvantage on attack rolls against the target. The target also can't be charmed, frightened, or possessed by them. If the target is already charmed, frightened, or possessed by such a creature, the target has advantage on any new saving throw against the relevant effect.
\end{document}
